\documentclass[11pt,a4paper]{article}

\usepackage[slovak]{babel}
\usepackage[utf8]{inputenc}
\usepackage[T1]{fontenc}
\usepackage{lmodern}
\usepackage[margin=2.5cm]{geometry}
\usepackage{amsmath,amssymb,amsthm}
\usepackage{hyperref}
\usepackage{microtype}
\usepackage{parskip}

\hypersetup{
    colorlinks=true,
    linkcolor=black,
    citecolor=black,
    urlcolor=blue
}

\title{Komunikácia s adaptívnym modelom prijímateľa:\\
prehľad konceptuálneho rámca}

\author{krse}

\date{\today}

\begin{document}

\maketitle

\begin{abstract}
\noindent
Klasická Shannonova teória informácie chápe prijímateľa ako pasívny
dekóder, ktorý rekonštruuje vyslanú správu s minimálnou distorziou.
Reálny prijímateľ však nie je optimálny: má obmedzený výpočtový výkon,
vlastný model sveta, cieľ rozhodovania a individuálnu históriu
interakcií. Tento článok navrhuje konceptuálny rámec, ktorý tieto
aspekty explicitne zahŕňa. Hlavnými komponentmi sú: (i) rozšírenie
rate-distortion funkcie o model a cieľ prijímateľa, (ii) učenie
parametrov prijímateľa zo spoločnej histórie komunikácie,
(iii) pravdepodobnostný rozhodovací model, ktorý meratelnou funkciou
vyjadruje \emph{cenu adaptácie} a porovnáva ju s očakávaným ziskom,
čím umožňuje racionálne rozhodnutie o tom, či je túto cenu vôbec
vhodné akceptovať, a (iv) termodynamické limity takejto adaptácie.
Ústrednou tézou je, že optimalizácia prenosu nie je zadarmo: každé
zlepšenie efektivity vyžaduje výpočet, pamäť, čas alebo energiu, a
niekedy táto cena prevýši hodnotu získanej informácie.  Cieľom je
pomenovať otvorené otázky a pripraviť terén pre formálne výsledky;
článok nedokazuje nové vety, ale identifikuje miesta, kde by mohli
vzniknúť.
\end{abstract}

\section{Úvod}

Shannonova rate-distortion teória~\cite{shannon1948,shannon1959}
stanovuje fundamentálnu medzu efektivity prenosu: zdroj $X$ sa
komprimuje do $R$ bitov a dekóder rekonštruuje $\hat{X}$ s chybou
$d(X,\hat{X}) \le D$. Krivka $R(D)$ je minimálny rate pri danej
distorzii a Shannon dokázal, že je dosiahnuteľná (s vhodne dlhým
kódovým blokom) nezávisle od prijímateľa.

Wyner a Ziv~\cite{wynerziv1976} rozšírili teóriu pre prípad, keď
prijímateľ disponuje vedľajšou informáciou $S$ so známym spoločným
rozdelením $P(X,S)$. Rate klesá na
\[
    R_{WZ}(D) \;=\; \min_{P(U|X)} [I(X;U) - I(S;U)]
\]
s rekonštrukciou $\hat{X} = \kappa(U,S)$.

Obe formulácie však predpokladajú, že prijímateľ je optimálny dekóder
s dokonalou znalosťou štatistík. V praxi to neplatí. Jazyk, dialóg
ľudí a strojov, sieťová komunikácia s rôznymi koncovými uzlami,
federované učenie, všetko sú systémy, kde prijímateľ je
\emph{agent} s vlastným modelom, cieľom a obmedzeniami. Tento článok
zhrnuje rámec, ktorý tieto prvky zavádza ako prvotriedne premenné.

\section{Model prijímateľa}

Prijímateľa charakterizujeme trojicou $(M, G, \kappa)$:
\begin{itemize}
    \item $M$: \emph{model sveta}, odhadované rozdelenie nad $X$,
          vo všeobecnosti odlišné od skutočného $P_X$.
    \item $G$: \emph{cieľ}, úloha s akčným priestorom $\mathcal{A}$
          a loss funkciou $\ell(a, X)$.
    \item $\kappa_M : \mathcal{I} \to \mathcal{A}$:
          \emph{rozhodovacia funkcia} parametrizovaná modelom $M$,
          nie ľubovoľná.
\end{itemize}

\noindent Navrhovaná rate-distortion funkcia:
\begin{equation}
    R_{M,G}(L) \;=\; \min_{f: \mathcal{X} \to \mathcal{I}}
    \mathbb{E}[|f(X)|]
    \quad \text{s.t.} \quad
    \mathbb{E}[\ell(\kappa_M(f(X)), X)] \le L.
    \label{eq:rmg}
\end{equation}

Odlišnosti oproti~$R_{WZ}$:
\begin{enumerate}
    \item distorzia meraná \emph{task-lossom} $\ell$ miesto~$d(X,\hat{X})$,
    \item $\kappa$ je \emph{fixovaná modelom $M$}, nie voľne optimalizovaná,
    \item očakávanie pod skutočným $P_X$, ale $\kappa$ používa
          \emph{vnímaný} model $M$, takže mismatch je vlastná zložka
          straty.
\end{enumerate}

Neformálna dekompozícia očakávanej straty:
\[
    \mathbb{E}[\ell] \;\approx\;
    \underbrace{L^\star(R)}_{\text{rate-limited}}
    + \underbrace{\Delta(M, P_X)}_{\text{model mismatch}}
    + \underbrace{\Delta(\kappa_M, \kappa^\star)}_{\text{bounded compute}}.
\]
Druhý a tretí člen nie sú odstrániteľné zvyšovaním~$R$, čo je zásadne
nové oproti Shannonovi.

\section{História ako zdroj adaptácie}

Parametre $(M, G, \kappa)$ vo všeobecnosti nepoznáme a v praxi sa
líšia medzi prijímateľmi. Jediný reálne dostupný zdroj informácie
o konkrétnom prijímateľovi je \emph{história} spoločnej komunikácie.
Analogicky s univerzálnym zdrojovým kódovaním~\cite{rissanen1978,
lempelziv1977} platí:
\[
    R_T \;=\; R^\star \;+\; \frac{c(k)}{T}
\]
kde $R^\star$ je optimálny rate pri plnej znalosti prijímateľa,
$T$~dĺžka histórie, $k$ zložitosť modelu. Pre veľké~$T$:
\begin{equation}
    \lim_{T \to \infty} \frac{1}{T}\sum_{t=1}^T R_t \;=\; h(X, A),
    \label{eq:entropy-rate}
\end{equation}
kde $h(X, A)$ je \emph{entropický rate} spoločného procesu zdroja
a reakcií prijímateľa. Toto je neredukovateľná dolná hranica.

Komunikácia sa teda prirodzene delí na tri fázy:
\begin{enumerate}
    \item \emph{Cold phase}: univerzálny kód, zber histórie.
    \item \emph{Warm phase}: postupná adaptácia na odhadnutý
          $(M, G, \kappa)$.
    \item \emph{Steady state}: $R_T \to R_{M,G}(L)$, penalty
          $c(k)/T$ mizne.
\end{enumerate}

V prípade obojsmernej interakcie~\cite{kaspi1985,orlitsky1990} je
konvergencia rýchlejšia, pretože každá výmena nesie dva signály
(obsah + nepriamy signál o cieli).

Ak máme k dispozícii históriu \emph{viacerých} prijímateľov, môžeme
odhadnúť \emph{populačný manifold} v priestore parametrov. Pre nového
prijímateľa potom stačí lokalizácia niekoľkými pozorovaniami
a adaptačná penalty klesá na $c(d)/T$, kde $d \ll k$ je vnútorná
dimenzia manifoldu (hierarchický Bayes, meta-learning).

\section{Rozhodovací rámec: kedy adaptovať}

Centrálna téza tohto článku znie: optimalizácia prenosu nie je
zadarmo. Adaptácia stojí čas~$\tau(T)$, pamäť, energiu a výpočet.
Preto potrebujeme \emph{meratelnú funkciu}, ktorá vyjadrí cenu
adaptácie a postaví ju proti očakávanému zisku, aby sme sa mohli
racionálne rozhodnúť, či túto cenu vôbec akceptovať. Bez takejto
funkcie sa optimalizácia stáva dogmou, pričom v niektorých prípadoch
je jej cena vyššia než hodnota získanej informácie.

Definujeme očakávaný zisk a očakávaný náklad:
\[
    \text{Benefit} = N_{\text{future}} \cdot [R_{\text{univ}} - R_{\text{adapt}}(T)] \cdot c_{\text{bit}},
\]
\[
    \text{Cost}(T) = C_{\text{comp}} \cdot \tau(T) + C_{\text{energy}} + C_{\text{mem}}.
\]

Deterministické rozhodovanie:
\[
    \text{adaptuj} \iff \text{Benefit} > \text{Cost}.
\]

Keďže všetky vstupné veličiny (gap, dĺžka vzťahu, $\tau$) sú neisté
a žiadni dvaja prijímatelia nie sú identickí, prirodzenejšia je
pravdepodobnostná verzia. Definujme:
\begin{equation}
    P_{\text{improve}}(k)
    \;=\; \Pr\!\big[\,G \cdot N \cdot c_{\text{bit}} > C_{\text{comp}} \cdot \tau
    \,\big|\, \text{obs}_{1:k}\,\big],
    \label{eq:pimprove}
\end{equation}
kde $G$ je gap $H(X) - h(X,A)$, $N$ očakávaný počet budúcich správ.
Funkcia $P_{\text{improve}}(k) \in [0,1]$ je \emph{meratelný výstup},
ktorý priamo podkladá rozhodnutie:
\[
    \text{adaptuj} \iff P_{\text{improve}}(k) > \theta.
\]
Prah~$\theta$ vyjadruje toleranciu rizika a ochotu platiť neistú cenu.

Pri Gaussovskom priore na populáciu $(G, N)$ má $P_{\text{improve}}$
uzavretú formu cez distribučnú funkciu~$\Phi$, teda je \emph{exaktne
vyčísliteľná} zo štatistík histórie a populácie. Presnosť odhadu
populačných parametrov rastie ako $1/\sqrt{N_{\text{pop}}}$.

Funkcia $P_{\text{improve}}(k)$ nevie odpovedať pre $k=0$ bez
populačného prioru, keďže rozhodovací model vždy vyžaduje buď
individuálnu históriu, alebo amortizovanú populačnú skúsenosť.
Samotný výpočet $P_{\text{improve}}$ má tiež cenu, ktorá musí byť
nižšia než priemerná úspora z rozhodnutí, ktoré generuje; inak sa
rozhodovací model stáva ďalším problémom, ktorý potrebuje vlastné
rozhodnutie (princíp \emph{selektívnej adaptácie}).

\section{Termodynamické limity}

Každá adaptácia modelu zahŕňa operácie nad pamäťou. Landauerov
princíp~\cite{landauer1961} kladie dolnú hranicu na energiu
nevratnej zmeny bitu:
\[
    E_{\min} = kT \ln 2.
\]
V praxi moderné systémy pracujú cca $10^{10}\times$ nad Landauerovou
hranicou, takže termodynamická zložka je zanedbateľná. V limitoch
(reverzibilné, kvantové systémy, vysoko hustá pamäť) sa stáva
dominantnou.

Globálne: druhý termodynamický zákon implikuje, že každá lokálna
redukcia informačnej entropie sa vykupuje väčším nárastom entropie
okolia. Informačná efektivita nie je zadarmo ani v tom najzákladnejšom
fyzikálnom zmysle, pretože sa vždy platí teplom odovzdaným do okolia.

\section{Hranice modelu}

Model je abstrakcia nad symbolovou a sémantickou vrstvou komunikácie;
fyzikálnu vrstvu (analógové veličiny, modulácia, šum kanála)
predpokladá za vyriešenú prostredníctvom klasického
\emph{channel coding theorem}. Pod pojmom \uv{dáta} rozumieme bity
prenesené kanálom; informácia, znalosť a múdrosť vznikajú až
u prijímateľa prostredníctvom interakcie dát s jeho modelom.

\emph{Úplná efektivita} prenosu nie je dosiahnuteľná z piatich
nezávislých dôvodov: (1)~šum kanála, (2)~neúplná znalosť prijímateľa,
(3)~konečný čas prenosu, (4)~konečný čas adaptácie,
(5)~ireducibilná entropia $h(X,A)$ spoločného procesu. Rámec sa
sústredí na druhý a štvrtý bod.

\section{Otvorené otázky}

\begin{itemize}
    \item \textbf{Gap teorém}: existuje presná horná/dolná hranica
          na $R_{M,G}(L) - R_{WZ}(L)$ ako funkcia $\text{KL}(P_X\|M)$?
    \item \textbf{Rate-distortion pre manifold prijímateľov}: aký je
          optimálny rate na reprezentáciu prijímateľa s danou
          kvantizačnou chybou? Táto otázka spája Shannonovu rate-distortion
          teóriu s manifold learningom a nonparametric Bayes.
    \item \textbf{Optimal stopping}: existuje uzavretá forma prahu
          $\theta^\star(\tau, N)$ pre pravdepodobnostné rozhodnutie~\eqref{eq:pimprove}?
    \item \textbf{Task-aware entropický rate}: pre akú triedu cieľov
          $G$ sa entropický rate $h(X,A)$ znižuje pod $h(X)$? Možno
          charakterizovať tieto triedy informačne-geometricky?
    \item \textbf{Ekvivalencia s information bottleneck}~\cite{tishby1999}:
          je kategorizácia prijímateľov ekvivalentná úlohe IB s task
          lossom ako relevanciou?
\end{itemize}

\section{Záver}

Prezentovaný rámec konsoliduje viaceré myšlienkové prúdy:
rate-distortion s vedľajšou informáciou, univerzálne zdrojové
kódovanie, semantickú komunikáciu~\cite{popovski2022}, meta-learning
a online adaptáciu. Žiaden z jednotlivých výsledkov nie je nový, nové
je ich \emph{integrované zarámovanie} s explicitným rozhodovacím
modulom a termodynamickým pohľadom. Ústredným motívom je konštrukcia
meratelnej funkcie pre cenu optimalizácie: bez nej sa adaptácia robí
slepo, s ňou sa stáva racionálnym rozhodnutím, ktoré občas znie
\uv{neadaptuj}, lebo vo svete nie je nič zadarmo. Cieľom ďalšej práce
je previesť otvorené otázky na formálne tvrdenia a dokázať aspoň jedno
netriviálne.

\bibliographystyle{plain}
\begin{thebibliography}{99}

\bibitem{shannon1948}
C.~E.~Shannon, \emph{A mathematical theory of communication},
Bell System Technical Journal, vol.~27, 1948.

\bibitem{shannon1959}
C.~E.~Shannon, \emph{Coding theorems for a discrete source with
a fidelity criterion}, IRE Nat.\ Conv.\ Rec., 1959.

\bibitem{wynerziv1976}
A.~D.~Wyner, J.~Ziv, \emph{The rate-distortion function for source
coding with side information at the decoder}, IEEE Trans.\ Inf.\ Theory,
vol.~22, 1976.

\bibitem{rissanen1978}
J.~Rissanen, \emph{Modeling by shortest data description},
Automatica, vol.~14, 1978.

\bibitem{lempelziv1977}
J.~Ziv, A.~Lempel, \emph{A universal algorithm for sequential data
compression}, IEEE Trans.\ Inf.\ Theory, vol.~23, 1977.

\bibitem{kaspi1985}
A.~H.~Kaspi, \emph{Two-way source coding with a fidelity criterion},
IEEE Trans.\ Inf.\ Theory, vol.~31, 1985.

\bibitem{orlitsky1990}
A.~Orlitsky, \emph{Worst-case interactive communication I: Two
messages are almost optimal}, IEEE Trans.\ Inf.\ Theory, vol.~36, 1990.

\bibitem{landauer1961}
R.~Landauer, \emph{Irreversibility and heat generation in the
computing process}, IBM J.\ Research and Development, vol.~5, 1961.

\bibitem{tishby1999}
N.~Tishby, F.~C.~Pereira, W.~Bialek, \emph{The information bottleneck
method}, Proc.\ 37th Allerton Conf., 1999.

\bibitem{popovski2022}
P.~Popovski, O.~Simeone, F.~Boccardi, D.~Gündüz, O.~Sahin,
\emph{Semantic-effectiveness filtering and control for post-5G
wireless connectivity}, Journal of the Indian Institute of Science,
2020; a následne literatúra o semantic communication, review 2022.

\end{thebibliography}

\end{document}
